% !PRINTED 544
\weberpage{544}
%<<CONT
samen Theiler, und folglich kann man nach 3. die Zahlen $\tau_1,
\tau_2, \ldots$ in $\mathfrak{o}$ so bestimmen, dass
\weq{8}{\varepsilon_1\tau_1 + \varepsilon_2\tau_2 + \cdots = 1}
wird. Aus (7) folgt aber, wenn man mit $\tau_1, \tau_2, \ldots$
multiplicirt und addirt, und dann
\[
\xi = \tau_1\xi_1 + \tau_2\xi_2 + \cdots, \quad
\eta = \tau_1\eta_1 + \tau_2\eta_2 + \cdots
\]
setzt, mittelst (8) die zu beweisende Gleichung (5).

\begin{center}
\S.\,150.

\gesperrt{Der Fermat'sche Satz.}
\end{center}

Aus der Theorie der Congruenzen lassen sich Folgerungen ziehen, die den
aus dem \gesperrt{Fermat'schen Lehrsatze} abgeleiteten Sätzen der
rationalen Zahlentheorie genau entsprechen, von denen hier die
wichtigsten, späterer Anwendung wegen, besprochen werden müssen.

Wir wollen unter $\pi$ ein Primfunctional $f^{\mathrm{ten}}$ Grades
verstehen, also, wenn $p$ die durch $\pi$ theilbare natürliche
Primzahl ist,
\weq{1}{N_a(\pi) = p^f}
setzen (\S.\,140). Ist dann $\alpha$ irgend eine durch $\pi$ nicht
theilbare Zahl in $\mathfrak{o}$, so wird, wie wir schon im vorigen
Paragraphen gesehen haben, das Product $\alpha\xi$ zugleich mit der
Zahl $\xi$ ein volles Restsystem nach dem Modul $\pi$ durchlaufen.
Lassen wir die durch $\pi$ theilbare Zahl weg, so bleiben $N_a(\pi) - 1$
Zahlen übrig, und wenn wir das Product bilden, so folgt
\weq{2}{\alpha^{p^f-1}\Pi(\xi) \equiv \Pi(\xi) \ (\mathrm{mod}\ \pi),}
wenn $\Pi(\xi)$ das Product aller Zahlen eines vollen Restsystems (mit
Ausschluss der Null) bedeutet, und daher durch $\pi$ nicht theilbar
ist. Demnach folgt aus (2)
\weq{3}{\alpha^{p^f-1} \equiv 1 \ (\mathrm{mod}\ \pi),}
oder, wenn man mit $\alpha$ multiplicirt,
\weq{4}{\alpha^{p^f} \equiv \alpha \ (\mathrm{mod}\ \pi),}
und in der letzten Form gilt der Satz auch noch, wenn $\alpha$ durch
$\pi$ theilbar ist.

Die Formeln (3), (4), die genau dem \gesperrt{Fermat'schen Lehrsatze}
entsprechen (Bd.\,I, \S.\,136), sollen auch hier als \gesperrt{Fermat'scher
Lehrsatz} bezeichnet werden.
